\makeatletter
\providecommand{\sgareaderlabel}[2]{%
  \begingroup
  \def\@currentlabel{#2}%
  \phantomsection\label{#1}%
  \endgroup}
\makeatother

\chapter{Application to the Fundamental Group}
\label{X}

\medskip
Throughout this expos\'e, $X$ will denote a locally noetherian prescheme,
$Y$ a closed subset of $X$, $U$ a variable open neighborhood of $Y$ in $X$,
and $\widehat X$ the formal completion of $X$ along $Y$
(\textup{EGA}~I, 10.8). For every prescheme $Z$, let $\Et(Z)$ denote the
category of \'etale coverings of $Z$, and let $\LLcat(Z)$\sgareaderlabel{pX.1}{1} denote the category
of coherent locally free modules on $Z$.

\section*{1. Comparison of $\Et(\widehat X)$ and $\Et(Y)$}
\sgareaderlabel{X.1}{1}

Let $\mathcal I$ be an ideal of definition of $Y$ in $X$. For every
$n\in\mathbb N$, set
\[
  Y_n=\left(Y,\left.(\mathcal O_X/\mathcal I^{\,n+1})\right|_Y\right).
\]
The $Y_n$ form a direct system of ordinary preschemes, or also of formal
preschemes when the structure sheaves are endowed with the discrete topology.
We know (EGA~I, 10.6.2) that $\widehat X$ is the direct limit, in the category
of formal preschemes, of the direct system $(Y_n)$. We also know (EGA~I,
10.13) that to give a formal $\widehat X$-prescheme $R$ of \emph{finite type}
is to give a direct system of ordinary $Y_n$-preschemes $R_n$ of
\emph{finite type}, such that
\[
  R_n\simeq R_{n+1}\times_{Y_{n+1}}Y_n.
\]
Moreover, for $R$ to be an étale covering of $\widehat X$, it is necessary
and sufficient that, for every $n$, $R_n$ be an étale covering of $Y_n$.

That said, it is easy to see that nilpotent elements make no difference to
étale coverings (SGA~1, 8.3); in other words, the base-change functor
\[
  \Et(Y_{n+1})\longrightarrow\Et(Y_n)
\]
is an equivalence of categories for every $n\in\mathbb N$. Hence:

\medskip
\sgareaderlabel{X.1.1}{1.1}\noindent\textbf{Proposition 1.1.}\quad
\emph{With the notation introduced above, the natural functor
\[
  \Et(\widehat X)\longrightarrow\Et(Y)
\]
is an equivalence of categories (cf.\ SGA~1, 8.4).}

\section*{2. Comparison of $\Et(Y)$ and $\Et(U)$, for varying $U$}
\sgareaderlabel{X.2}{2}

We shall introduce two conditions from which the comparison theorem announced
above will follow easily. Let $X$ be a locally noetherian prescheme, and let
$Y$ be a closed subset of $X$. We say that the pair $(X,Y)$ satisfies the
\emph{Lefschetz condition}, written $\Lefc(X,Y)$,\sgareaderlabel{pX.2}{2} if, for every open subset
$U$ of $X$ containing $Y$ and every locally free coherent sheaf $E$ on $U$,
the natural homomorphism
\[
  \Gamma(U,E)\longrightarrow\Gamma(\widehat X,\widehat E)
\]
is an isomorphism.

We say that the pair $(X,Y)$ satisfies the \emph{effective Lefschetz
condition}, written $\Leffc(X,Y)$, if $\Lefc(X,Y)$ holds and, in addition,
for every locally free coherent sheaf $\mathcal E$ on $\widehat X$, there
exist an open neighborhood $U$ of $Y$, a locally free coherent sheaf $E$ on
$U$, and an isomorphism
\[
  \widehat E\simeq\mathcal E.
\]

These conditions hold in two important examples:

\medskip
\sgareaderlabel{X.2.1}{2.1}\noindent\textbf{\textit{Example 2.1}}\footnote{\sgareaderlabel{noteX.2.1}{*}\textit{Editor's note.}
One can slightly improve (i): see Mme Raynaud (Raynaud~M., ``Théorèmes de
Lefschetz en cohomologie des faisceaux cohérents et en cohomologie étale.
Application au groupe fondamental,'' \emph{Ann. Sci. Éc. Norm. Sup. (4)}
\textbf{7} (1974), pp.~29--52, Corollaries I.1.4 and I.5). Condition (ii) can
be improved so as to dispense with the depth conditions along $Y$ (see
Theorem~3.3 of \emph{loc. cit.} for a precise statement). The proof of this
latter point is very technical; moreover, the article cited above gives only
indications of proof and refers to an earlier, detailed version published in
the \emph{Bulletin de la Société mathématique de France}.}.\ ---\ 
Let $A$ be a noetherian ring, and let $t\in\mathfrak r(A)$ be an
$A$-regular element belonging to the radical $\mathfrak r(A)$ of $A$.
Assume that $A$ is a quotient of a regular local ring and that $A$ is
complete for the $t$-adic topology (for example, that $A$ is complete for
the $\mathfrak r(A)$-adic topology). Set
\[
  X'=\operatorname{Spec}(A),\qquad Y'=V(t),
\]
and, further, set
\[
  x=\mathfrak r(A),\qquad X=X'\setminus\{x\},\qquad
  Y=Y'\setminus\{x\}.
\]
Thus $X$ is open in $X'$ and $Y=X\cap Y'$. Then:

\begin{enumerate}[label=(\roman*),leftmargin=2.5em]
\item If, for every prime ideal $\mathfrak p$ of $A$ such that
$\dim A/\mathfrak p=1$ (i.e. for every closed point of $X$), we have
$\depthop A_{\mathfrak p}\geq 2$, then $\Lefc(X,Y)$ holds;

\item if, moreover, for every prime ideal $\mathfrak p$ of $A$ such that
$t\in\mathfrak p$ and $\dim A/\mathfrak p=1$ (i.e. for every closed point
of $Y$), we have $\depthop A_{\mathfrak p}\geq 3$, then $\Leffc(X,Y)$
holds.
\end{enumerate}

Let us first prove that, for every open neighborhood $U$ of $Y$ in $X$, the
complement of $U$ in $X$ is a union of finitely many closed points (in $X$).
Observe that $U$ is open in $X$, hence in $X'$, and therefore
$Z'=X'\setminus U$ is closed. Let $I$ be an ideal defining $Z'$; it is
enough to prove that $A/I$ has dimension $1$. Since
\[
  Z'\cap Y'=\{x\},
\]
the ring $A/(I+(t))$ is Artinian, and the conclusion follows from the
``Hauptidealsatz.''

The \emph{first hypothesis is equivalent to}: ``for every prime ideal
$\mathfrak p$ of $A$ with $\mathfrak p\ne\mathfrak r(A)$, one has
\[
  \operatorname{depth} A_{\mathfrak p}
  \geq 3-\dim(A/\mathfrak p).
\]
Indeed, $A$ is a quotient of a regular ring, so we may apply \hyperref[VIII.2.3]{VIII, 2.3} to
the prescheme $X'$, the closed subset $\{x\}$, and the coherent sheaf
$\mathcal O_{X'}$, observing that
$c(\mathfrak p)=\dim(A/\mathfrak p)$ for
$\mathfrak p\in X=X'\setminus\{x\}$ (since $x$ is the closed point of
$X'$).

\medskip
\noindent\colorbox{noteblue}{\parbox{0.95\linewidth}{\footnotesize
\textbf{Source oddity SGA2-X-L3357-U-VS-X-SRCDEF-001.}
The French authority prints $\mathfrak p\in U=X'-\{x\}$, although $U$ was
just quantified as an arbitrary open neighborhood and line 3345 defines
$X=X'-\{x\}$. The English therefore uses $X$ here. The French file remains
unchanged; the same-edition PDF confirms only the printed $U$, not the
repair.}}

Let $U$ be an open neighborhood of $Y$ in $X$, and let $E$ be a locally
free $\mathcal O_U$-module. Set $Z=X\setminus U$, and let
$u\colon U\to X$ be the canonical immersion. We shall first prove that
$u_*(E)$ is a coherent $\mathcal O_X$-module, or, equivalently, that
$\mathcal H^i_Z(E')$ is coherent for $i=0,1$, where $E'$ is a coherent
extension of $E$ to $X$. For this, we apply Theorem \hyperref[VIII.2.1]{VIII, 2.1} to the
prescheme $X$, the closed subset $Z$, and the coherent sheaf $E'$. It is
enough to verify that, for every point $p\in U$ such that $c(p)=1$, one has
$\operatorname{depth}E'_p\geq1$, where we have set
\[
  c(p)=\operatorname{codim}
  \bigl(\overline{\{p\}}\cap Z,\overline{\{p\}}\bigr).
\]
Now suppose $p\in U$ and $c(p)=1$, and let $\mathfrak p$ denote the ideal
of $A$ corresponding to $p$. Then $\dim(A/\mathfrak p)=2$, since the
complement of $U$ is a union of finitely many closed points and $A$ is a
quotient of a regular ring. Moreover, $E$ is locally free, so for every
$p\in\operatorname{Supp}E$ one has
$\operatorname{depth}E_p=\operatorname{depth}\mathcal O_{U,p}$. Finally,
if $p\in U$ and $c(p)=1$, then
\[
  \operatorname{depth}E'_p
  =\operatorname{depth}E_p
  =\operatorname{depth}\mathcal O_{U,p}
  =\operatorname{depth}A_{\mathfrak p}
  \geq3-2=1.
\]

We must now prove that the natural homomorphism
\begin{equation}\label{eq:X.1}
  \Gamma(U,E)\longrightarrow\Gamma(\widehat X,\widehat E)
\end{equation}
is an isomorphism. Set $\overline{\!E}=u_*(E)$. The sheaf
$\overline{\!E}$ is coherent and has depth at least $2$ at every closed
point of $X$. It follows that $R^if_*(\overline{\!E})$ is coherent for
$i=0,1$, where $f\colon X\to X'$ denotes the canonical immersion of $X$
into $X'=\operatorname{Spec}(A)$ (Exposé VIII). Applying \hyperref[IX.1.5]{IX, 1.5}, we
conclude that
\begin{equation}\label{eq:X.2}
  \Gamma(U,\overline{\!E})
  \longrightarrow
  \Gamma(\widehat X,\widehat{\overline{\!E}})
\end{equation}
is an isomorphism, since $A$ is complete for the $t$-adic topology.

There is a commutative diagram
\[
\xymatrix@C=18mm@R=9mm{
  \Gamma(X,\overline{\!E})
    \ar[rr]^{\simeq}
    \ar[dr]_{\simeq}
  && \Gamma(U,E)\ar[dl] \\
  & \Gamma(\widehat X,\widehat E) &
}
\]
and the conclusion follows.

\medskip
Now let $\mathcal E$ be a locally free coherent sheaf on $\widehat X$. If
one has proved that $\mathcal E$ is algebraizable, that is, isomorphic to the
formal completion of a coherent $\mathcal O_X$-module $E$, it is easy to see
that $E$ is locally free in a neighborhood of $Y$, and hence to prove
$\Leffc(X,Y)$. Let $\widehat{X'}$ be the formal spectrum of $A$ for the
$t$-adic topology; it is identified with the formal completion of $X'$ along
$Y'$. Denote by $f$ the canonical immersion of $X$ into $X'$, by $f'$ the
canonical immersion of $Y$ into $Y'$, and by $\widehat f$ the morphism
obtained by passing to the completions. For $\mathcal E$ to be algebraizable,
it suffices that $\widehat f_*(\mathcal E)$ be a coherent
$\mathcal O_{\widehat{X'}}$-module, since $A$ is complete for the $t$-adic
topology.
\sourceoddity{SGA2-X-L3384-PUSHFORWARD-BASE-SRCDEF-001}{The French authority
prints $\mathcal O_{\widehat X}$ as the coefficient ring of the pushforward.
Because $\widehat f\colon\widehat X\to\widehat{X'}$, the target uses
$\mathcal O_{\widehat{X'}}$. Only this coefficient ring is corrected.}
Set $\mathcal I=t\mathcal O_{\widehat X}$; this is an ideal of definition of
$\widehat X$.

For every $n\geq0$, set
\[
  \mathcal E_n=\mathcal E/\mathcal I^{n+1}\mathcal E.
\]
At every closed point $y\in Y$, the depth of $\mathcal E_0$ is at least two;
indeed, $t$ is an $A$-regular element, and therefore
\[
  \depthop\mathcal O_{Y,y}
  =\depthop\mathcal O_{X,y}-1\geq2.
\]
\sourceoddity{SGA2-X-L3386-Y0-VS-Y-SRCDEF-001}{The French authority prints
$\mathcal O_{Y_0,y}$, but $Y_0$ is undefined in this argument. Since
$Y=V(t)\cap X$, the regular-element quotient gives $\mathcal O_{Y,y}$.}
It follows that $\widehat f_*(\mathcal E)$ is coherent (\hyperref[IX.2.3]{IX~2.3)}.
\hfill\textsc{q.e.d.}

\medskip
\sgareaderlabel{X.2.2}{2.2}\noindent\textbf{\textit{Example 2.2}}\quad
\textup{(This will make it possible to compare the fundamental group of a
projective variety with that of a hyperplane section).}

Let $K$ be a field, and let $X$ be a proper $K$-prescheme. Let $\mathcal L$
be an ample invertible $\mathcal O_X$-module. Let
$t\in\Gamma(X,\mathcal L)$ be an $\mathcal O_X$-regular element; this means
that, for every open subset $U$ and every isomorphism
\[
  u\colon \mathcal L|_U \xrightarrow{\ \sim\ } \mathcal O_U,
\]
the section $u(t)$ is a non-zero-divisor in $\mathcal O_U$ (a condition that
does not depend on $u$). Let $Y=V(t)$ be the subscheme of $X$ defined by the
equation $t=0$\footnote{\textit{Editor's note.} Condition (ii) is
superfluous; see the editor's note to \hyperref[X.2.1]{Example~2.1} (original running p.~90).}.
Then:

\begin{enumerate}[label=(\roman*),leftmargin=2.5em]
\item If, for every closed point $x$ of $X$, one has
$\depthop\mathcal O_{X,x}\geq2$, then $\Lefc(X,Y)$ holds;

\item if, moreover, for every closed point $y\in Y$, one has
$\depthop\mathcal O_{X,y}\geq3$, then $\Leffc(X,Y)$ holds.
\end{enumerate}

\medskip
This example will be treated in detail in Exposé~XII.

\medskip
Let $S$ be a prescheme. It is known (EGA~II, 6.1.2) that the functor which
assigns to every finite flat covering $r\colon R\to S$ the
$\mathcal O_S$-algebra $r_*(\mathcal O_R)$ induces an equivalence between
the category of finite flat coverings of $S$ and the category of coherent
locally free $\mathcal O_S$-algebras.
\sourceoddity{SGA2-X-L3402-OS-VS-OX-SRCDEF-001}{The French authority prints
$\mathcal O_X$ in both coefficient-ring positions. Since $S$ is arbitrary
and $r\colon R\to S$, the direct image and the equivalent algebra category
are over $\mathcal O_S$.}
Let $U$ be an open neighborhood of $Y$, and let $r\colon R\to U$ be a
finite flat covering of $U$. Let $\widehat R$ be the finite flat covering of
$\widehat X$ obtained from it by base change. Then
\[
  \widehat r_*(\mathcal O_{\widehat R})
  \simeq \widehat{r_*(\mathcal O_R)}.
\]

Assume now that $\Lefc(X,Y)$ holds. This implies that, for every $U$, the
inverse-image functor
\[
  \mathsf L(U)\longrightarrow\mathsf L(\widehat X)
\]
is fully faithful. Indeed, let $E$ and $F$ be two coherent locally free
$\mathcal O_U$-modules; $\cHom(E,F)$ is likewise coherent and locally free.
\sourceoddity{SGA2-X-L3408-DUPLICATE-HOM-SRCDEF-001}{The French authority
prints the preceding sentence twice consecutively. The target states the
identical assertion once.}
By hypothesis, the natural map
\[
  \Gamma\bigl(U,\cHom(E,F)\bigr)
  \longrightarrow
  \Gamma\bigl(\widehat X,\widehat{\cHom(E,F)}\bigr)
\]
is an isomorphism, which proves the assertion, since $\cHom$ commutes with
formal completion when all the modules involved are locally free. Formal
completion also commutes with tensor products. It follows that the functor
which assigns to every coherent locally free $\mathcal O_U$-algebra
$\mathcal A$ the $\mathcal O_{\widehat X}$-algebra $\widehat{\mathcal A}$
is fully faithful. More precisely, if $E$ is a coherent locally free
$\mathcal O_U$-module, commutative $\mathcal O_U$-algebra structures on $E$
correspond bijectively to commutative $\mathcal O_{\widehat X}$-algebra
structures on $\widehat E$.
\sourceoddity{SGA2-X-L3412-BIJECTION-MISSING-FIRST-SIDE-SRCDEF-001}{The
French authority introduces a bijective correspondence but prints only its
$\mathcal O_{\widehat X}$ side. The target restores the omitted
$\mathcal O_U$-algebra structures on $E$, as required by the preceding
fully faithful tensor-compatible completion functor.}

\medskip
\sgareaderlabel{X.2.3}{2.3}\noindent\textbf{\textit{Proposition 2.3.}}
Let $X$ be a locally noetherian prescheme, and let $Y$ be a closed subset of
$X$. Let $\widehat X$ be the formal completion of $X$ along $Y$. For every
open subset $U$ of $X$ containing $Y$, let $L_U$ (respectively $P_U$,
respectively $E_U$) denote the functor that assigns to every locally free
coherent $\mathcal O_U$-module (respectively every finite flat covering of
$U$, respectively every \'{e}tale covering of $U$) its pullback along
$\widehat X\to X$.

\begin{enumerate}[label=(\roman*),leftmargin=2.5em]
\item If $\Lefc(X,Y)$ holds, then, for every open neighborhood $U$ of $Y$,
the functors $L_U$, $P_U$, and $E_U$ are fully faithful.

\item If $\Leffc(X,Y)$ holds, then, for every locally free coherent
$\mathcal O_{\widehat X}$-module $\mathcal E$, there are an open
neighborhood $U$ of $Y$ and a locally free coherent $\mathcal O_U$-module
$E$ such that
\[
  L_U(E)\simeq\mathcal E.
\]
The analogous assertions hold for finite flat coverings and for \'{e}tale
coverings, with $P_U$ and $E_U$, respectively.
\end{enumerate}

\medskip
\noindent\textit{Proof.}
Part~(i) has already been proved.

Part~(ii) follows from (i) and the hypothesis, at least for $L_U$ and $P_U$.
Moreover, if $R$ is an \'{e}tale covering of $\widehat X$, there are an open
neighborhood $U$ of $Y$ in $X$ and a finite flat covering $R'$ of $U$ such
that
\[
  \widehat{R'}\simeq R.
\]
By restriction this gives a covering $R''$ of $Y$, which is \'{e}tale by
\hyperref[X.1.1]{Proposition~1.1}. Hence the restriction of $R'$ to some neighborhood $U'$ of
$Y$ is \'{e}tale.
\hfill$\square$

\medskip
\sgareaderlabel{X.2.4}{2.4}\noindent\textbf{\textit{Corollary 2.4.}}
Suppose that $\Lefc(X,Y)$ holds. A finite flat covering $R$ of an open
neighborhood $U$ of $Y$ is connected if and only if
$R\times_U\widehat X$ is connected. In particular, $Y$ is connected if and
only if $U$ is connected, and this is also equivalent to $X$ being
connected.

\medskip
\noindent\textit{Proof.}
Indeed, a locally ringed space $(X,\mathcal O_X)$ is connected if and only if
$\Gamma(X,\mathcal O_X)$ is not a direct product of two nonzero rings. But,
by $\Lefc(X,Y)$, one has
\[
  \Gamma\bigl(U,r_*\mathcal O_R\bigr)
  \simeq
  \Gamma\bigl(\widehat X,\widehat r_*\mathcal O_{\widehat R}\bigr).
\]

\medskip
\sgareaderlabel{X.2.5}{2.5}\noindent\textbf{\textit{Corollary 2.5.}}
Suppose that $\Lefc(X,Y)$ holds. Then, for every open neighborhood $U$ of
$Y$, the functor
\[
  \Etc(U)\longrightarrow\Etc(Y)
\]
is fully faithful. Suppose that $\Leffc(X,Y)$ holds. Then, for every
\'{e}tale covering $R$ of $Y$, there are an open neighborhood $U$ of $Y$
and a covering $R'$ of $U$ such that
$R'\times_U Y\simeq R$.

\medskip
\sgareaderlabel{X.2.6}{2.6}\noindent\textbf{\textit{Corollary 2.6\textsuperscript{(3)}.}}
Suppose that $\Lefc(X,Y)$ holds and that $Y$ is connected. Then every open
neighborhood $U$ of $Y$ is connected, and the natural homomorphism
$\pi_1(Y)\longrightarrow\pi_1(U)$ is surjective. Suppose, moreover, that
$\Leffc(X,Y)$ holds. Then the natural homomorphism
\[
  \pi_1(Y)\longrightarrow\varprojlim_{U}\,\pi_1(U)
\]
is an isomorphism. (N.B. A ``base point'' in $Y$ is understood to have been
chosen and is also taken as a base point in $X$ for the definition of the
fundamental groups.)

\medskip
\noindent All of this follows immediately from \hyperref[X.1.1]{Propositions~1.1} and~\hyperref[X.2.3]{2.3}.

\section*{Editor's note (3)}
Combined with \hyperref[X.3.3]{Proposition~3.3} and criteria \hyperref[XII.2.4]{XII.2.4} and \hyperref[XII.3.4]{XII.3.4}, this yields the
following relative Lefschetz theorem. Let $f\colon X\to S$ be a projective and
flat morphism of connected noetherian schemes, and let $D$ be an effective
relative Cartier divisor on $X$ that is relatively ample. If, for every
$s\in S$, the depth of $X_s$ at every closed point is $\geq 2$, then $D$ is
connected and, for every open subset $U$ of $X$ containing $D$, the map
\[
  i_U\colon \pi_1(D)\longrightarrow\pi_1(U)
\]
is surjective. If, in addition, the depth of $X_s$ at every closed point of
$D_s$ is $\geq 3$ and the local rings of $X$ at its closed points are
pure---for example, are complete intersections (cf. \hyperref[X.3.4]{X, 3.4)}---then $i_X$ is an
isomorphism.

Cf. J.-B. Bost, ``Lefschetz theorem for Arithmetic Surfaces,''
\emph{Ann. Sci. \'Ec. Norm. Sup. (4)} \textbf{32} (1999), pp.~241--312,
Theorems~1.1 and~2.1.

When $X$ is merely a smooth projective surface that is geometrically connected
over a field, $D$ is always connected and
\[
  \pi_1(D)\longrightarrow\pi_1(U)
\]
is surjective (where $U$ is an open subset containing $D$), provided only that
$D$ is nef and has self-intersection $>0$; cf. \emph{loc. cit.}, Theorem~2.3,
and also Theorem~2.4 for surfaces assumed only normal and complete.

In the case of an arithmetic surface (normal and quasi-projective) $X$ over a
ring of integers $\mathcal O_K$, Bost, improving results of Y.~Ihara
(Y.~Ihara, ``Horizontal divisors on arithmetic surfaces associated with
Bely\u{\i} uniformizations,'' in \emph{The Grothendieck theory of dessins
d'enfants (Luminy, 1993)}, London Math. Soc. Lect. Note Series, vol.~200,
Cambridge Univ. Press, Cambridge, 1994, pp.~245--254, or \emph{loc. cit.},
Corollary~7.2), proved that if a point
$P\in X(\mathcal O_K)$, which plays the role of the divisor $D$ in the
geometric situation, satisfies certain positivity conditions, then the map
\[
  \pi_1(X)\longrightarrow\pi_1(\operatorname{Spec}\mathcal O_K)
\]
induced by the projection is an isomorphism whose inverse is the map
\[
  \pi_1(\operatorname{Spec}\mathcal O_K)\longrightarrow\pi_1(X)
\]
induced by $P$; cf. \emph{loc. cit.}, Theorem~1.2.

\medskip
\noindent{\footnotesize\emph{Accepted source-normalization note
(SGA2-X-L3456-AU-DESSUS-NORMALIZATION-001@1).} The corrected French control
selects \textup{au-dessus} from
\texttt{\string\sisi\{au dessus\}\{au-dessus\}}; the target renders this as
``over.'' This is an explicit source normalization, not a source defect.}

\section*{3. Comparison of $\pi_1(X)$ and $\pi_1(U)$}
\sgareaderlabel{X.3}{3}

\medskip
\sgareaderlabel{X.3.1}{3.1}\noindent\textbf{Definition 3.1.}\quad
Let $X$ be a prescheme and $Z$ a closed subset of $X$. Set
$U=X\setminus Z$. We say that the pair $(X,Z)$ is \emph{pure} if, for every
open subset $V$ of $X$, the functor
\[
\begin{aligned}
  \Et(V)&\longrightarrow\Et(V\cap U),\\
  V'&\longmapsto V'\times_V(V\cap U)
\end{aligned}
\]
is an equivalence of categories\renewcommand{\thefootnote}{*}\footnote{For a
notion that is more satisfactory in certain respects, cf. the commentary
\hyperref[XIV.1.6]{XIV, 1.6(d)}.}.

\medskip
\noindent{\footnotesize\emph{Accepted source-normalization note
(SGA2-X-L3464-MAPSTO-NORMALIZATION-001@1).} The corrected French control
selects \texttt{\string\mto} from
\texttt{\string\sisi\{\string\rightsquigarrow\}\{\string\mto\}}; the target
renders the object map with $\longmapsto$. This is an explicit source
normalization, not a source defect.}

\medskip
\sgareaderlabel{X.3.2}{3.2}\noindent\textbf{Definition 3.2.}\quad
Let $A$ be a noetherian local ring. Set $X=\Spec A$. Let $\rr(A)$ be the
radical of $A$, and let $x=\rr(A)$ be the closed point of $X$. We say that
$A$ is \emph{pure} if the pair $(X,\{x\})$ is pure.

\medskip
\noindent We leave to the reader the task of not proving the following
proposition:

\medskip
\sgareaderlabel{X.3.3}{3.3}\noindent\textbf{Proposition 3.3.}\quad
Let $X$ be a locally noetherian prescheme and let $Z$ be a closed subset of
$X$. The pair $(X,Z)$ is pure if and only if, for every $z\in Z$, the local
ring $\Oo_{X,z}$ is pure\renewcommand{\thefootnote}{**}\footnote{Compare the
noncommutative case of \hyperref[XIV.1.8]{XIV, 1.8}, whose proof is essentially the same as that
of \hyperref[X.3.3]{Proposition~3.3}.}.

\medskip
\noindent{\footnotesize\color{notebluetext}\emph{Source-defect note
(SGA2-X-L3477-FERME-VS-FERMEE-SRCDEF-001).} French line~3477 reads
\emph{une partie ferm\'e de $X$}; the masculine adjective \emph{ferm\'e}
disagrees with the feminine noun \emph{partie}. The target renders the
intended sense as ``a closed subset of $X$.'' The admitted French control is
unchanged.}

\medskip
\noindent That said, the following theorem is the main result of this
section:

\medskip
\sgareaderlabel{X.3.4}{3.4}\noindent\textbf{Theorem 3.4 (Purity theorem)}%
\renewcommand{\thefootnote}{(4)}\footnote{\textsc{N.D.E.}: For the history of
the methods employed, see Grothendieck's letter of October~1, 1961 to Serre,
\emph{Correspondance Grothendieck-Serre}, edited by Pierre Colmez and
Jean-Pierre Serre, Documents Math\'ematiques, vol.~2, Soci\'et\'e
Math\'ematique de France, Paris, 2001.}\textbf{.}

\begin{enumerate}
  \renewcommand{\labelenumi}{(\roman{enumi})}
  \item A regular noetherian local ring of dimension $\geq 2$ is pure
  (Zariski--Nagata purity theorem).
  \item A noetherian local ring of dimension $\geq 3$ that is a complete
  intersection is pure.
\end{enumerate}

\medskip
\noindent Recall that a local ring $A$ is said to be a \emph{complete
intersection} if and only if there exist a \emph{regular} noetherian local
ring $B$ and a $B$-regular sequence $(t_1,\ldots,t_k)$ of elements of the
radical $\mathfrak r(B)$ of $B$ such that
\[
  A \simeq B/(t_1,\ldots,t_k).
\]

\medskip
\noindent In this connection, let us note that it would be less ambiguous to
say that $A$ is an \emph{absolute} complete intersection, in contrast to the
situation already encountered, in which $X$ is a locally noetherian prescheme
(not necessarily regular) and $Y$ is a closed subset of $X$ that is said to
be ``locally a set-theoretic complete intersection \emph{in $X$}.''

\medskip
\noindent Let us first prove a few lemmas.

\medskip
{\itshape
\sgareaderlabel{X.3.5}{3.5}\noindent\textbf{Lemma 3.5.} --- Let $X$ be a locally noetherian prescheme,
and let $U$ be an open subset of $X$. Set $Z=X\setminus U$. Let
$i\colon U\to X$ be the canonical immersion of $U$ into $X$. The following
conditions are equivalent:
\begin{enumerate}[label=\textup{(\roman*)},leftmargin=2.2em,itemsep=0.45em,topsep=0.45em]
\item For every open subset $V$ of $X$, if $V'=V\cap U$, the restriction
functor $F\longmapsto F|_{V'}$ from the category of locally free coherent
$\mathcal O_V$-modules to the category of locally free coherent
$\mathcal O_{V'}$-modules is fully faithful;
\item the natural homomorphism $\mathcal O_X\to i_*\mathcal O_U$ is an
isomorphism;
\item for every $z\in Z$, one has $\prof\mathcal O_{X,z}\geq 2$.
\end{enumerate}
}

We have already seen (\Ref{III}~\Ref{III.3.3}) the equivalence of (ii)
and (iii). Let us show that (ii) implies (i). Let $F$ and $G$ be two
locally free coherent $\Oo_V$-Modules. Then $\SheafHom(F,G)$ is also
locally free and coherent, so
$$
\SheafHom(F,G)\longrightarrow
i_\ast\bigl(\SheafHom(\sisi{F|{|V'}}{F_{|V'}},
\sisi{G|{|V'}}{G_{|V'}})\bigr)
$$
is an isomorphism, and therefore
$$
\Hom(F,G)\simeq
\Hom(\sisi{F|{|V'}}{F_{|V'}},\sisi{G|{|V'}}{G_{|V'}}).
$$
Conversely, take $F=G=\OX$ and apply (i) to every open subset $V$ of
$X$.

\medskip
\noindent Here is a useful ``descent lemma'':

\medskip
{\itshape
\sgareaderlabel{X.3.6}{3.6}\noindent\textbf{Lemma 3.6.} --- Let $X$ be a locally noetherian prescheme,
and let $Z$ be a closed subset of $X$. Set $U=X\setminus Z$. Suppose that
the homomorphism $\mathcal O_X\to i_*\mathcal O_U$ is an isomorphism. Let
$f\colon X_1\to X$ be a faithfully flat and quasi-compact morphism. Set
$Z_1=f^{-1}(Z)$. If the pair $(X_1,Z_1)$ is pure, then so is $(X,Z)$.
}

\medskip
\noindent\textit{Proof of \hyperref[X.3.6]{Lemma 3.6}.}\quad
Observe that the hypothesis
$\Oo_X\simeq i_*\Oo_U$ is preserved under flat base extension, since $i$ is
quasi-compact and, in this case, formation of the direct image commutes with
inverse image. This hypothesis implies that the functor
\[
  \Et(V)\longrightarrow\Et(U\cap V)
\]
defined by
\[
  V'\longmapsto V'\times_V(V\cap U)
\]
is fully faithful, as follows from the interpretation of an étale covering
in terms of a coherent locally free algebra. It remains to prove effectivity.
One may, for example, form the fiber square $X_2$ and the fiber cube $X_3$ of
$X_1$ over $X$, and observe that a faithfully flat and quasi-compact
morphism is a morphism of universal effective descent for the fibered
category of étale coverings over the category of preschemes. The conclusion
then follows formally\renewcommand{\thefootnote}{(*)}\footnote{Cf. J.~Giraud,
\emph{Méthode de la descente}, Mémoire no.~2 du \emph{Bulletin de la Société
Mathématique de France} (1964).}.

\medskip
{\itshape
\sgareaderlabel{X.3.7}{3.7}\noindent\textbf{Remark 3.7.} --- We proved along the way that if
$\mathcal O_X\to i_*\mathcal O_U$ is an isomorphism, then $X$ is connected if
and only if $U$ is connected, and then $\pi_1(U)\to\pi_1(X)$ is surjective.
}

\medskip
{\itshape
\sgareaderlabel{X.3.8}{3.8}\noindent\textbf{Corollary 3.8.} --- Let $A$ be a noetherian local ring.
Assume that $\depthop A\geq 2$. If $\widehat{A}$ is pure, then $A$ is pure.
}

\medskip
\noindent\emph{Proof.}\quad This follows from \hyperref[X.3.5]{Lemma~3.5} and \hyperref[X.3.6]{Lemma~3.6}.

\medskip
\noindent The following lemma is the essential point in the proof of the purity theorem:

\medskip
{\itshape
\sgareaderlabel{X.3.9}{3.9}\noindent\textbf{Lemma 3.9.} --- Let $A$ be a noetherian local ring, and
let $t\in\mathfrak r(A)$ be an $A$-regular element. Suppose that $A$ is
complete for the $t$-adic topology and is moreover a quotient of a regular
local ring (for example, if $A$ is complete). Set $B=A/tA$.
\begin{enumerate}[label=\textup{(\roman*)},leftmargin=2.2em,itemsep=0.6em,topsep=0.5em]
\item If, for every prime ideal $\mathfrak p$ of $A$ such that
$\dim A/\mathfrak p=1$, one has $\depthop A_{\mathfrak p}\geq 2$, then
the purity of $B$ implies the purity of $A$.

\item If, for every prime ideal $\mathfrak p$ of $A$ such that
$\dim A/\mathfrak p=1$, one has $\depthop A_{\mathfrak p}\geq 2$, if
$A_{\mathfrak p}$ is pure whenever $t\notin\mathfrak p$, and
if\renewcommand{\thefootnote}{(5)}\footnote{\textsc{N.D.E.}: This last
condition can be improved; cf. the editor's note~(1) on p.~90.}
$\depthop A_{\mathfrak p}\geq 3$ whenever $t\in\mathfrak p$, then the
purity of $A$ implies the purity of $B$.
\sourceoddity{SGA2-X-L3551-MISSING-EST-SRCDEF-001}{The French authority
prints \emph{si $A_{\mathfrak p}$ pur lorsque $t\notin\mathfrak p$},
omitting the finite copula \emph{est}. The English supplies ``is'' in the
unambiguous intended condition. The French authority remains unchanged.}
\end{enumerate}
}

\medskip
Let $X'=\operatorname{Spec}(A)$ and let $Y'=V(t)$, which we identify with
the spectrum of $B$. Let $x=\mathfrak r(A)$, and set
$X=X'\setminus\{x\}$ and $Y=Y'\setminus\{x\}=X\cap Y'$. Denote by
$\widehat{X'}$ the formal spectrum of $A$ for the $t$-adic topology; it is
identified with the formal completion of $X'$ along $Y'$.

Since $A$ is complete for the $t$-adic topology, observe that
$\Et(X')\longrightarrow\Et(\widehat{X'})$ is an equivalence of categories.
Likewise, $\Et(\widehat{X'})\longrightarrow\Et(Y')$ is an equivalence by
\hyperref[X.1.1]{Proposition~1.1}; hence $\Et(X')\longrightarrow\Et(Y')$ is an equivalence of
categories.

We prove \textup{(i)}. Consider the diagrams
\[
\xymatrix@C=12mm@R=8mm{
  X' & X\ar[l] &&
  \Et(X')\ar[r]^{a}\ar[d]_{c} & \Et(X)\ar[d]^{b}\\
  Y'\ar[u] & Y\ar[l]\ar[u] &&
  \Et(Y')\ar[r]^{d} & \Et(Y)
}
\]
We have just seen that $c$ is an equivalence; $d$ is also an equivalence by
the hypothesis that $B$ is pure; and finally $b$ is fully faithful, as was
seen in \hyperref[X.2.1]{Example~2.1}; cf. \hyperref[X.2.3]{Proposition~2.3}\textup{(i)}.

We prove \textup{(ii)}. This time suppose that $A$ is pure, so $a$ is an
equivalence; likewise $c$. Let us show that $b$ is an equivalence. By
\hyperref[X.2.1]{Example~2.1}, we know that $\Leffc(X,Y)$ holds, so $b$ is already fully
faithful; let us prove that it is essentially surjective. Apply
\hyperref[X.2.3]{Proposition~2.3}\textup{(ii)}, noting that if $U$ is an open neighborhood of
$Y$ in $X$, then $X\setminus U$ is a union of finitely many closed points.
The pair $(X,X\setminus U)$ is therefore pure by \hyperref[X.3.3]{Proposition~3.3}, since at
any such point $\mathfrak p$ one has
$\mathcal O_{X,\mathfrak p}=A_{\mathfrak p}$, which is pure by hypothesis.
This proves the assertion.

\medskip
\textit{Proof of the purity theorem.}
We first prove \textup{(i)} by induction on the dimension. Let $A$ be a
noetherian local ring of dimension~$2$. Set
$X'=\Spec(A)$, $x=\mathfrak r(A)$, and $X=X'\setminus\{x\}$. We have
$\depthop A=2$. We may therefore apply \hyperref[X.3.5]{Lemma~3.5} to the pair
$(X',\{x\})$, and hence
$\Et(X')\longrightarrow\Et(X)$ is fully faithful.

Now let $r\colon R\longrightarrow X$ be an \'{e}tale covering defined by
the coherent, locally free, \'{e}tale $\mathcal O_X$-algebra
$\mathcal A=r_*(\mathcal O_R)$. Let $i\colon X\longrightarrow X'$ denote
the canonical immersion. I claim that
$i_*(\mathcal A)=\mathcal B$ is a coherent $\mathcal O_{X'}$-algebra.
Indeed, it suffices to apply the ``finiteness theorem,'' \hyperref[VIII.2.3]{VIII, 2.3}. I also
claim that this algebra has depth at least~$2$ at $x$: it is the direct image
of an $\mathcal O_X$-module, with $X=X'\setminus\{x\}$. Since $A$ is a
regular ring of dimension~$2$, we have
\[
  \pdop \mathcal B+\depthop \mathcal B=\dim A=2,
\]
where $\pdop \mathcal B$ denotes the projective dimension of $\mathcal B$.
Thus $\pdop \mathcal B=0$, so $\mathcal B$ is projective and therefore
free. Consequently, $\mathcal B$ defines a finite flat covering of
$X'=\Spec(A)$. The set of points of $X'$ at which this covering is not
\'{e}tale is a closed subset of $X'$ whose defining ideal is principal,
namely the discriminant ideal of $\mathcal B/A$. By construction this closed
subset is contained in $x=\mathfrak r(A)$; it is therefore empty because
$\dim A=2$.

Let $A$ now be a regular noetherian local ring with $\dim A=n\geq 3$.
Suppose that \textup{(i)} has been proved for rings of dimension less than
$n$. To prove that $A$ is pure, we may assume by \hyperref[X.3.8]{Corollary~3.8} that $A$ is
complete. Choose $t\in\mathfrak r(A)$ whose image in
$\mathfrak r(A)/\mathfrak r(A)^2$ is nonzero. Then $B=A/tA$ is a regular
noetherian local ring of dimension $n-1$, and is therefore pure because
$n-1\geq 2$. The conclusion follows from \hyperref[X.3.9]{Lemma~3.9}\textup{(i)}, which is
applicable because $A$ is complete.

\medskip
\textit{Let us prove \textup{(ii)}.}
Let $A$ be a noetherian local ring of dimension at least~$3$. Suppose that
there exist a regular noetherian local ring $B$ and a $B$-sequence
$(t_1,\ldots,t_k)$ such that
\[
  A\simeq B/(t_1,\ldots,t_k).
\]
We prove that $A$ is pure by induction on $k$. If $k=0$, this follows from
\textup{(i)}. Suppose that $k\geq1$ and that the result has been proved for
$k'<k$. By \hyperref[X.3.8]{Corollary~3.8}, we may assume that $A$, and hence also $B$, is
complete. Set
\[
  C=B/(t_1,\ldots,t_{k-1}).
\]
Then $A\simeq C/t_kC$, and $t_k$ is $C$-regular. By the induction hypothesis,
$C$ is pure; it therefore suffices to verify that \hyperref[X.3.9]{Lemma~3.9}\textup{(ii)} is
applicable. In the notation of that lemma, its $A$ and $B$ are here $C$ and
$A$, respectively. We have $\dim C\geq4$; hence, for every prime ideal
$\mathfrak p$ of $C$ such that $\dim(C/\mathfrak p)=1$, one has
$\depthop C_{\mathfrak p}\geq3$. Moreover, $C_{\mathfrak p}$ is a complete
intersection with $k'\leq k-1$, and is therefore pure by the induction
hypothesis.

\medskip
{\itshape
\sgareaderlabel{X.3.10}{3.10}\noindent\textbf{Theorem 3.10.} --- Let $X$ be a locally noetherian
prescheme, and let $Y$ be a closed subset of $X$. Suppose that
$\Leff(X,Y)$ holds (cf. \hyperref[X.2.1]{Examples~2.1} and~\hyperref[X.2.2]{2.2}). Suppose further that, for
every open neighborhood $U$ of $Y$ and every $x\in X\setminus U$, the
local ring $\mathcal O_{X,x}$ is either regular of dimension $\geq 2$ or
a complete intersection of dimension $\geq 3$. Then
\[
  \pi_0(Y)\longrightarrow\pi_0(X)
\]
is a bijection, and, if $X$ is connected,
\[
  \pi_1(Y)\longrightarrow\pi_1(X)
\]
is an isomorphism.
}

\medskip

There is nothing more to prove. Observe that in the two cited
\hyperref[X.2.1]{Examples~2.1} and~\hyperref[X.2.2]{2.2}, the complement of $U$ is a union of finitely many
closed points; it follows that the hypothesis on the dimension of
$\mathcal O_{X,x}$ is not farcical.
